%======================================================================
\section{Numerical experiments}\label{sec:experiments}
%======================================================================

Our numerical experiments evaluate four properties of the proposed method: (i) stop-gradient interface targets eliminate the bias of joint training (\cref{sec:exp-interface}, \cref{thm:stationary}); (ii) the error scales with the learning grid, as predicted by the error budget (\cref{sec:exp-budget}, \cref{thm:attainable}); (iii) local log-signatures capture the path information of path-dependent and non-Markovian problems (\cref{sec:exp-information,sec:exp-rough}, \cref{lem:information}); and (iv) every trained model returns a primal--dual certificate that compares favorably with established methods (\cref{sec:exp-benchmarks}, \cref{thm:primaldual}).

%----------------------------------------------------------------------
\subsection{Test problems and protocol}\label{sec:exp-setup}
%----------------------------------------------------------------------

\Cref{tab:problems} lists the contracts. Assets follow correlated geometric Brownian motions with $S_0=100$, $r=5\pct$, and $\sigma=20\pct$; \cref{sec:exp-rough} adds a rough Bergomi model. Reference values come from binomial and average trees, closed forms, and the primal--dual values of \cite{andersen2004}.

\begin{table}[htbp]
\centering\footnotesize
\caption{Test problems. $K$ is the number of BSDE layers.}
\label{tab:problems}
\setlength{\tabcolsep}{4pt}
\begin{tabular}{@{}l l c c c l@{}}
\toprule
Contract & Payoff at $T_k$ & $d$ & $K$ & $T$ & Reference\\
\midrule
Geometric Asian call & $\left(\exp(\tfrac1T\!\int_0^T\log S_t\dd t)-100\right)^+$ & 1 & 1 & 1 & closed form\\
Bermudan put & $(\kappa-S_{T_k})^+$, \ $\kappa\in\{90,100,110\}$ & 1 & 4--24 & 1 & binomial tree\\
Amerasian put & $\left(100-\tfrac1{T_k}\!\int_0^{T_k}S_t\dd t\right)^+$ & 1 & 4--12 & 1 & average tree\\
Bermudan basket put & $(100-\tfrac1d\sum_iS^i_{T_k})^+$, \ $\rho=0.3$ & 1--10 & 4 & 1 & --\\
Geometric basket call & $\left((\prod_iS^i_T)^{1/d}-100\right)^+$ & 1--10 & 1 & 1 & closed form\\
Bermudan max-call & $(\max_iS^i_{T_k}-100)^+$, \ $q=10\pct$, $\rho=0$ & 2, 5 & 9 & 3 & \cite{andersen2004}\\
Rough Bergomi put & $(\kappa-S_{T_k})^+$, \ $\kappa\in\{90,100,110\}$ & 1 & 4, 12 & 1 & --\\
\bottomrule
\end{tabular}
\end{table}

The encoder computes level-three local log-signatures of the time-augmented path (level two for $d\ge5$) and drives a Neural RDE with $64$ hidden units shared by all layers; we refer to this model as LogSig--NRDE. The value and hedge heads are two-layer networks of width $64$. Unless stated otherwise, the learning grid has $N=24$ steps, and paths are simulated on about $1000$ fine steps, exactly for geometric Brownian motion and by the hybrid scheme \cite{bennedsen2017} for rough Bergomi. Training uses Adam with batch size $512$ and $2000$ updates and detaches the interface targets unless stated otherwise. Prices and hedges are evaluated on $4\times10^4$ common test paths, $L_0$ on $2\times10^5$ further paths with a control variate, and $U_0$ with the network hedge as the dual martingale. Results are means over three seeds, and comparisons between methods use paired differences with $95\pct$ hierarchical bootstrap intervals.

%----------------------------------------------------------------------
\subsection{Interface targets}\label{sec:exp-interface}
%----------------------------------------------------------------------

The four variants of \cref{fig:interface} share one forward computation and differ only in the gradient path of the interface residual. Under stop-gradient training, each value head becomes the conditional expectation of its target, and the stationary price equals $V_0^{\mathrm f}$ (\cref{thm:stationary}(i), \cref{lem:information}(iv)). By contrast, joint training adds a bias term driven by the preceding residual on the continuation region (\cref{thm:stationary}(ii)).


To assess hedging quality alongside pricing accuracy, \cref{tab:interface} also reports the $L^2$ hedge error $\|Z-\hZ\|_{L^2}$ and the excess hedging error, defined as the root-mean-square discounted profit and loss of the delta hedge rebalanced on the learning grid and held until the network's exercise date, minus the same quantity for the exact hedge on that grid. The excess hedging error thus isolates the part of the replication error attributable to the network. Consistent with \cref{thm:stationary}, stop-gradient training prices all contracts within $0.31\pct$ and reduces the hedge error of joint training by a factor of three at virtually identical cost ($412$\,s versus $405$\,s). EMA targets hedge equally well but misprice the Bermudan puts by $0.72$--$1.91\pct$. Layerwise training yields comparable prices with a larger hedge error, and its full-budget version requires nearly twice the training time.

\begin{table}[htbp]
\centering\footnotesize
\caption{Training variants ($K=4$). Relative price error $|y_0-V_0|/V_0$ in \%; hedge error $\|Z-\hZ\|_{L^2}$, excess hedging error, and training time for the Bermudan put with $\kappa=100$. Bold: best in column; orange: main weakness of an alternative.}
\label{tab:interface}
\setlength{\tabcolsep}{4pt}
\begin{tabular}{@{}l cccc c c c@{}}
\toprule
& \multicolumn{4}{c}{price error [\%]} & hedge error & excess hedging & time\\
\cmidrule(lr){2-5}
Variant & $\kappa=90$ & $\kappa=100$ & $\kappa=110$ & Amerasian & [$\times10^{-3}$] & error & [s]\\
\midrule
\weak{joint} \cite{huang2026} & \weak{5.99} & \weak{3.53} & \weak{1.36} & \weak{5.65} & \weak{1.36} & \weak{0.052} & \weak{405}\\
\rowours{\textbf{stop-gradient (ours)}} & \rowours{\textbf{0.15}} & \rowours{\textbf{0.08}} & \rowours{0.07} & \rowours{0.31} & \rowours{0.42} & \rowours{0.007} & \rowours{412}\\
EMA target                   & \weak{0.72} & \weak{1.34} & \weak{1.91} & \textbf{0.14} & \textbf{0.28} & \textbf{0.006} & 507\\
layerwise                    & 0.17 & 0.15 & \textbf{0.03} & 0.89 & \weak{2.33} & \weak{0.042} & 188\\
layerwise, full budget       & 0.20 & 0.18 & 0.05 & 0.35 & 0.56 & 0.008 & \weak{759}\\
\bottomrule
\end{tabular}
\end{table}

The joint bias behaves as a first-order discretization effect: it halves with the learning step and persists under longer training (\cref{fig:bias}), as \cref{rem:stationary} predicts. Stop-gradient training removes it on every grid and converges to the reference value as training proceeds.

\begin{figure}[htbp]
\centering
\resizebox{\textwidth}{!}{% Joint vs stop-gradient: y0 bias against the learning step h and against
% the training budget (Bermudan put, K = 4, strike 100).  data/a2_*.dat
\begin{tikzpicture}
\begin{groupplot}[group style={group size=2 by 1, horizontal sep=1.9cm},
                  paper, width=6.9cm, height=5.1cm]
\nextgroupplot[
  title={(a) Bias against the learning step},
  xmode=log, ymode=log, log basis x=2,
  xlabel={learning step $h=T/N$}, ylabel={$(y_0-V_0)/V_0$ \ [\%]},
  xmin=0.008, xmax=0.11, ymin=0.002, ymax=20,
  xtick={0.0104167,0.0208333,0.0416667,0.0833333},
  xticklabels={$\tfrac1{96}$,$\tfrac1{48}$,$\tfrac1{24}$,$\tfrac1{12}$},
  ytick={0.01,0.1,1,10}, yticklabels={$0.01$,$0.1$,$1$,$10$},
  legend pos=south east]
\addplot[sGuide, domain=0.0095:0.1] {78.69*x};
\addplot[sJoint] table[x=h, y=joint_bias] {data/a2_nscan.dat};
\addplot[sStop]  table[x=h, y=stopgrad_bias] {data/a2_nscan.dat};
\node[font=\scriptsize, text=cInk!70, rotate=33] at (axis cs:0.017,2.2) {slope $1$};
\legend{,joint, stop-gradient}

\nextgroupplot[
  title={(b) Bias against the training budget},
  xmode=log, log basis x=2,
  xlabel={parameter updates}, ylabel={$(y_0-V_0)/V_0$ \ [\%]},
  xmin=800, xmax=10000, ymin=-0.6, ymax=4.4,
  xtick={1000,2000,4000,8000}, xticklabels={$1000$,$2000$,$4000$,$8000$},
  ytick={0,1,2,3,4},
  legend style={at={(0.97,0.58)}, anchor=east}]
\addplot[draw=none, fill=cBand, forget plot] coordinates {(800,-0.15) (10000,-0.15) (10000,0.15) (800,0.15)} \closedcycle;
\addplot[sJoint] table[x=steps, y=joint_bias] {data/a2_steps.dat};
\addplot[sStop]  table[x=steps, y=stopgrad_bias] {data/a2_steps.dat};
\draw[cInk!40, line width=0.5pt] (axis cs:800,0) -- (axis cs:10000,0);
\legend{joint, stop-gradient}
\end{groupplot}
\end{tikzpicture}
}
\caption{Price bias of joint and stop-gradient training, Bermudan put ($\kappa=100$, $K=4$). (a) Against the learning step; the dashed line has slope one. (b) Against the number of updates at $N=24$; the band marks $\pm0.15\pct$.}
\label{fig:bias}
\end{figure}

The advantage persists for many exercise dates (\cref{fig:kscan}a,b). Replacing the Bermudan maximum by identity links, however, removes the joint bias (\cref{fig:kscan}c), which shows that the bias originates in the exercise interface, in line with the continuation indicator of \cref{thm:stationary}(ii). The identity-link scan further shows that the price error of trained networks stays far inside the factor $K$ of \cref{thm:aposteriori}.

\begin{figure}[htbp]
\centering
\resizebox{\textwidth}{!}{% Relative price error against the number of layers K, three learning steps
% per layer.  data/cv_*.dat (mean and standard deviation over three seeds)
\begin{tikzpicture}
\begin{groupplot}[group style={group size=3 by 1, horizontal sep=1.25cm},
                  paper, width=5.25cm, height=5.0cm,
                  xmode=log, ymode=log, ymin=0.008, ymax=25,
                  ytick={0.01,0.1,1,10}, yticklabels={$0.01$,$0.1$,$1$,$10$},
                  xlabel={layers $K$}]
\nextgroupplot[title={(a) Bermudan put}, ylabel={$|y_0-V_0|/V_0$ \ [\%]},
               xmin=3, xmax=32, xtick={4,12,24}, xticklabels={$4$,$12$,$24$}]
\addplot[sJoint] plot[error bars/.cd, y dir=both, y explicit] table[x=K, y=joint, y error plus=joint_sd, y error minus=joint_lo] {data/cv_kscan_berm.dat};
\addplot[sLayer] plot[error bars/.cd, y dir=both, y explicit] table[x=K, y=layerwise, y error plus=layerwise_sd, y error minus=layerwise_lo] {data/cv_kscan_berm.dat};
\addplot[sStop]  plot[error bars/.cd, y dir=both, y explicit] table[x=K, y=stopgrad, y error plus=stopgrad_sd, y error minus=stopgrad_lo] {data/cv_kscan_berm.dat};

\nextgroupplot[title={(b) Amerasian put},
               xmin=3, xmax=16, xtick={4,12}, xticklabels={$4$,$12$}]
\addplot[sJoint] plot[error bars/.cd, y dir=both, y explicit] table[x=K, y=joint, y error plus=joint_sd, y error minus=joint_lo] {data/cv_kscan_amer.dat};
\addplot[sLayer] plot[error bars/.cd, y dir=both, y explicit] table[x=K, y=layerwise, y error plus=layerwise_sd, y error minus=layerwise_lo] {data/cv_kscan_amer.dat};
\addplot[sStop]  plot[error bars/.cd, y dir=both, y explicit] table[x=K, y=stopgrad, y error plus=stopgrad_sd, y error minus=stopgrad_lo] {data/cv_kscan_amer.dat};

\nextgroupplot[title={(c) Identity links},
               xmin=0.8, xmax=16, xtick={1,2,4,12}, xticklabels={$1$,$2$,$4$,$12$},
               legend style={at={(0.97,0.97)}, anchor=north east}]
\addplot[sJoint] plot[error bars/.cd, y dir=both, y explicit] table[x=K, y=joint, y error plus=joint_sd, y error minus=joint_lo] {data/cv_identity.dat};
\addplot[sLayer] plot[error bars/.cd, y dir=both, y explicit] table[x=K, y=layerwise, y error plus=layerwise_sd, y error minus=layerwise_lo] {data/cv_identity.dat};
\addplot[sStop]  plot[error bars/.cd, y dir=both, y explicit] table[x=K, y=stopgrad, y error plus=stopgrad_sd, y error minus=stopgrad_lo] {data/cv_identity.dat};
\legend{joint, layerwise, stop-gradient}
\end{groupplot}
\end{tikzpicture}
}
\caption{Relative price error against the number of layers, three learning steps per layer; whiskers mark one standard deviation.}
\label{fig:kscan}
\end{figure}

%----------------------------------------------------------------------
\subsection{Error budget}\label{sec:exp-budget}
%----------------------------------------------------------------------

\Cref{thm:attainable} splits the error into learning-grid, monitoring, truncation, and network terms. The endpoint loss follows the grid floor $\varepsilon^\pi(Z^{\mathrm f})=O(h)$, and the hedge error decreases with it (\cref{fig:floor}a,b). At a fixed number of steps per layer, $K\cL(\theta)$ is nearly constant in $K$ (\cref{fig:floor}c), so the loss scales like the lower end $\varepsilon^\pi(Z^{\mathrm f})/K$ of \cref{thm:attainable}(a). By contrast, varying the truncation level from $M=1$ to $4$ (\cref{fig:floor}d), which increases the network size from about $29\,000$ to $54\,000$ parameters, keeps the endpoint loss at $1.28$ and the hedge error between $1.35\times10^{-3}$ and $1.54\times10^{-3}$, and refining the fine grid from $252$ to $2016$ steps keeps the loss within $1.26$--$1.29$ and the price error below $0.1\pct$. Consistent with \cref{lem:information,thm:attainable}, once $M$ captures the payoff state and the fine grid resolves the running integral, accuracy is governed by the learning grid, at a computational cost linear in $N$.

\begin{figure}[htbp]
\centering
\resizebox{\textwidth}{!}{% Loss floor and hedge error against h (geometric Asian call, K = 1) and the
% scaled floor K*L against K at three steps per layer (Bermudan put);
% (d) loss against the truncation level M and the fine grid N_f.
\begin{tikzpicture}
\begin{groupplot}[group style={group size=2 by 2, horizontal sep=1.6cm, vertical sep=1.7cm},
                  paper, width=6.6cm, height=4.6cm]
\nextgroupplot[title={(a) Endpoint loss}, xmode=log, ymode=log, log basis x=2,
  xlabel={$h$}, ylabel={$\mathcal L(\theta)$},
  xmin=0.008, xmax=0.11, ymin=0.2, ymax=4,
  xtick={0.0104167,0.0208333,0.0416667,0.0833333},
  xticklabels={$\tfrac1{96}$,$\tfrac1{48}$,$\tfrac1{24}$,$\tfrac1{12}$},
  ytick={0.25,0.5,1,2}, yticklabels={$0.25$,$0.5$,$1$,$2$}]
\addplot[sGuide, domain=0.0095:0.1] {31.09*x};
\addplot[sStop] table[x=h, y=J] {data/cv_asian_N.dat};
\node[font=\scriptsize, text=cInk!70, rotate=36] at (axis cs:0.0165,0.72) {slope $1$};

\nextgroupplot[title={(b) Hedge error}, xmode=log, ymode=log, log basis x=2,
  xlabel={$h$}, ylabel={$\|Z-\widehat Z\|_{L^2}$},
  xmin=0.008, xmax=0.11, ymin=1e-4, ymax=1e-2,
  xtick={0.0104167,0.0208333,0.0416667,0.0833333},
  xticklabels={$\tfrac1{96}$,$\tfrac1{48}$,$\tfrac1{24}$,$\tfrac1{12}$}]
\addplot[sStop] table[x=h, y=z] {data/cv_asian_N.dat};

\nextgroupplot[title={(c) Scaled floor}, xmode=log,
  xlabel={layers $K$}, ylabel={$K\,\mathcal L(\theta)$},
  xmin=3, xmax=32, ymin=0, ymax=24, xtick={4,12,24}, xticklabels={$4$,$12$,$24$},
  legend pos=south east]
\addplot[sJoint] table[x=K, y expr=\thisrow{K}*\thisrow{joint_J}] {data/cv_kscan_berm.dat};
\addplot[sStop]  table[x=K, y expr=\thisrow{K}*\thisrow{stopgrad_J}] {data/cv_kscan_berm.dat};
\legend{joint, stop-gradient}

\nextgroupplot[title={(d) Truncation and monitoring},
  xlabel={$M$ \ / \ fine steps $N_{\mathrm f}$}, ylabel={$\mathcal L(\theta)$},
  xmin=0.6, xmax=4.4, ymin=1.0, ymax=1.6, xtick={1,2,3,4},
  xticklabels={$1$/$252$,$2$/$504$,$3$/$1008$,$4$/$2016$},
  ytick={1.0,1.2,1.4,1.6}, legend pos=north east]
\addplot[sStop] table[x=M, y=J] {data/cv_asian_M.dat};
\addplot[color=cEma, mark=diamond*, mark size=2.6pt] table[x expr=ln(\thisrow{Nf}/126)/ln(2), y=J] {data/cv_asian_Nf.dat};
\legend{vary $M$, vary $N_{\mathrm f}$}
\end{groupplot}
\end{tikzpicture}
}
\caption{Error budget. (a), (b) Geometric Asian call against the learning step; the dashed line has slope one. (c) Bermudan put, three learning steps per layer. (d) Geometric Asian call against the truncation level $M$ ($N_{\mathrm f}=1008$) and the number of fine steps $N_{\mathrm f}$ ($M=3$).}
\label{fig:floor}
\end{figure}

%----------------------------------------------------------------------
\subsection{Path information}\label{sec:exp-information}
%----------------------------------------------------------------------

\Cref{thm:wellposed}(ii) represents each layer value as a functional of the price path. The Amerasian put makes this dependence explicit: a spot-state encoder $(t,S_t)$ as in \cite{huang2026} has no access to the running average and fails even when its parameter count is matched to that of our model (\cref{fig:encoders}a). The shared LogSig--NRDE encoder matches the same network supplied with the true running average and outperforms LSTM, signature-head, and backward deep signature \cite{bayraktar2024} encoders.

\begin{figure}[htbp]
\centering
\resizebox{\textwidth}{!}{% Path encoders, K = 4.  (a) Amerasian price error; (b)-(d) paired difference of the
% squared hedging error (encoder minus ours) with 95% interval; (e), (f) size and time
% on the Amerasian put.  data/bl_encoders.dat and the paired bootstrap output.
\begin{tikzpicture}
\pgfplotsset{forest/.style={paper, ymin=-0.6, ymax=8.6, y dir=reverse, ytick={0,...,8},
    ymajorgrids=false, xmajorgrids=true, axis lines*=box,
    xticklabel style={/pgf/number format/fixed}, scaled x ticks=false, clip=false},
  ci/.style={only marks, mark=*, mark size=1.7pt, color=cInk!80, mark options={fill=cInk!80},
    error bars/x dir=both, error bars/x explicit,
    error bars/error bar style={line width=0.8pt, color=cInk!80},
    error bars/error mark options={rotate=90, mark size=1.8pt, color=cInk!80}},
}
\begin{groupplot}[group style={group size=3 by 2, horizontal sep=0.45cm, vertical sep=1.85cm,
    yticklabels at=edge left},
  forest, width=5.0cm, height=5.6cm, yticklabels={{\textbf{\textcolor{cStop!80!black}{shared LogSig--NRDE (ours)}}},{separate encoders},{LogSig--NRDE + running average},{full-path LogSig--NRDE},{signature heads},{backward deep signature},{LSTM on LogSig},{Markov state, param.-matched},{Markov state $(t,S_t)$}},
  yticklabel style={font=\scriptsize, text=cInk}]
\nextgroupplot[title={(a) Price error}, xmode=log, xmin=0.08, xmax=300,
  xtick={0.1,1,10}, xticklabels={$0.1$,$1$,$10$}, xlabel={$|y_0-V_0|/V_0$ [\%]},
  xbar, bar width=6pt, bar shift=0pt, log origin=infty,
  nodes near coords, point meta=rawx, nodes near coords style={font=\tiny, text=cInk, anchor=west,
  xshift=0pt, /pgf/number format/.cd, fixed, precision=2}]
\addplot[fill=cStop, draw=cStop!75!black] coordinates {(0.3149,0)};
\addplot[fill=cInk!25, draw=cInk!50] coordinates {(0.1289,1) (0.3735,2) (0.3926,3) (0.4728,4) (0.4888,5) (0.5918,6)};
\addplot[fill=cJoint!65, draw=cJoint] coordinates {(18.252,7) (18.212,8)};
\nextgroupplot[title={(b) Bermudan put}, xmin=-0.03, xmax=0.19, xlabel={$\Delta$ sq.\ hedging error}]
\fill[cStop!9] (axis cs:0,\pgfkeysvalueof{/pgfplots/ymin}) rectangle (axis cs:\pgfkeysvalueof{/pgfplots/xmax},\pgfkeysvalueof{/pgfplots/ymax});
\draw[cStop!85!black, line width=1.1pt] (axis cs:0,\pgfkeysvalueof{/pgfplots/ymin}) -- (axis cs:0,\pgfkeysvalueof{/pgfplots/ymax});
\addplot[only marks, mark=diamond*, mark size=3pt, color=cStop!85!black] coordinates {(0,0)};
\addplot[ci] coordinates {(0.022,1) += (0.0140,0) -= (0.0220,0) (0.012,3) += (0.0050,0) -= (0.0070,0) (0.04,4) += (0.0060,0) -= (0.0060,0) (0.027,5) += (0.0060,0) -= (0.0060,0) (0.157,6) += (0.0200,0) -= (0.0210,0) (0.04,7) += (0.0080,0) -= (0.0070,0) (0.087,8) += (0.0090,0) -= (0.0090,0)};
\node[font=\tiny, text=cInk, anchor=north east, inner sep=2pt] at (axis cs:0,-0.55) {\textcolor{cStop!80!black}{ours}};
\nextgroupplot[title={(c) Basket put, $d=5$}, xmin=-0.012, xmax=0.05, xlabel={$\Delta$ sq.\ hedging error}]
\fill[cStop!9] (axis cs:0,\pgfkeysvalueof{/pgfplots/ymin}) rectangle (axis cs:\pgfkeysvalueof{/pgfplots/xmax},\pgfkeysvalueof{/pgfplots/ymax});
\draw[cStop!85!black, line width=1.1pt] (axis cs:0,\pgfkeysvalueof{/pgfplots/ymin}) -- (axis cs:0,\pgfkeysvalueof{/pgfplots/ymax});
\addplot[only marks, mark=diamond*, mark size=3pt, color=cStop!85!black] coordinates {(0,0)};
\addplot[ci] coordinates {(0.002,1) += (0.0020,0) -= (0.0020,0) (-0.002,3) += (0.0040,0) -= (0.0030,0) (0.019,4) += (0.0030,0) -= (0.0020,0) (0.014,5) += (0.0030,0) -= (0.0020,0) (0.034,6) += (0.0080,0) -= (0.0070,0) (0.03,7) += (0.0020,0) -= (0.0030,0) (0.04,8) += (0.0030,0) -= (0.0030,0)};
\node[font=\tiny, text=cInk, anchor=north east, inner sep=2pt] at (axis cs:0,-0.55) {\textcolor{cStop!80!black}{ours}};
\nextgroupplot[title={(d) Amerasian put}, xmin=-0.04, xmax=0.12, xlabel={$\Delta$ sq.\ hedging error}]
\fill[cStop!9] (axis cs:0,\pgfkeysvalueof{/pgfplots/ymin}) rectangle (axis cs:\pgfkeysvalueof{/pgfplots/xmax},\pgfkeysvalueof{/pgfplots/ymax});
\draw[cStop!85!black, line width=1.1pt] (axis cs:0,\pgfkeysvalueof{/pgfplots/ymin}) -- (axis cs:0,\pgfkeysvalueof{/pgfplots/ymax});
\addplot[only marks, mark=diamond*, mark size=3pt, color=cStop!85!black] coordinates {(0,0)};
\addplot[ci] coordinates {(-0.02,1) += (0.0070,0) -= (0.0060,0) (-0.009,3) += (0.0040,0) -= (0.0030,0) (0.074,4) += (0.0050,0) -= (0.0050,0) (0.048,5) += (0.0050,0) -= (0.0060,0) (0.046,6) += (0.0220,0) -= (0.0310,0)};
\draw[->, cJoint, line width=0.9pt] (axis cs:0.075,7) -- (axis cs:0.117,7) node[pos=0.45, above, font=\tiny, text=cInk, inner sep=1pt] {2.30};
\draw[->, cJoint, line width=0.9pt] (axis cs:0.075,8) -- (axis cs:0.117,8) node[pos=0.45, above, font=\tiny, text=cInk, inner sep=1pt] {2.31};
\node[font=\tiny, text=cInk, anchor=north east, inner sep=2pt] at (axis cs:0,-0.55) {\textcolor{cStop!80!black}{ours}};
\nextgroupplot[title={(e) Parameters}, xmin=0, xmax=260, xlabel={parameters [$\times10^{3}$]},
  xbar, bar width=6pt, bar shift=0pt, nodes near coords, point meta=rawx,
  nodes near coords style={font=\tiny, text=cInk, anchor=west, /pgf/number format/.cd, fixed, precision=0}]
\addplot[fill=cStop, draw=cStop!75!black] coordinates {(93.001,0)};
\addplot[fill=cInk!25, draw=cInk!50] coordinates {(193.225,1) (130.953,2) (93.001,3) (37.193,4) (37.194,5) (86.217,6) (93.347,7) (30.921,8)};
\nextgroupplot[title={(f) Training time}, xmin=0, xmax=1900, xticklabel style={/pgf/number format/1000 sep={}}, xlabel={time [s]},
  xbar, bar width=6pt, bar shift=0pt, nodes near coords, point meta=rawx,
  nodes near coords style={font=\tiny, text=cInk, anchor=west, /pgf/number format/.cd, fixed, precision=0, 1000 sep={}}]
\addplot[fill=cStop, draw=cStop!75!black] coordinates {(402,0)};
\addplot[fill=cInk!25, draw=cInk!50] coordinates {(1384,1) (435,2) (465,3) (92,4) (281,5) (93,6) (65,7) (62,8)};
\end{groupplot}
\end{tikzpicture}
}
\caption{Path encoders ($K=4$). (a) Relative price error on the Amerasian put, logarithmic scale. (b)--(d) Squared hedging error minus that of the shared LogSig--NRDE, with $95\pct$ interval; the shaded side favors our model. (e), (f) Model size and training time on the Amerasian put.}
\label{fig:encoders}
\end{figure}

Path information also improves hedging for Markovian problems (\cref{fig:encoders}b--d). All alternative encoders yield larger hedging errors on the Bermudan put, and on the five-asset basket put only the full-path LogSig--NRDE matches our shared local encoder. Moreover, sharing one encoder across layers is computationally decisive: compared with layer-specific encoders, it uses less than half the parameters and cuts training time by a factor of $3.4$ on the Amerasian put (\cref{fig:encoders}e,f), needs $162\,649$ instead of $463\,321$ parameters on the basket put, and delivers nearly identical primal--dual bounds.


%----------------------------------------------------------------------
\subsection{Rough volatility}\label{sec:exp-rough}
%----------------------------------------------------------------------

The rough Bergomi model \cite{bayer2016} with $H=0.07$, $\eta=1.9$, $\rho=-0.9$, and $\xi_0=0.09$ violates \cref{ass:obs}, since the price is no longer a Markov state and the spot variance is unobserved. \Cref{fig:rough} compares encoders and regression methods for the Bermudan put. Using only the price path, the LogSig--NRDE model attains tighter dual bounds than the Markov-state encoder and signature Longstaff--Schwartz, with equal or better lower bounds. Adding the spot variance as an input channel tightens the certificate further. Even then, the LogSig--NRDE encoder improves on the Markov state augmented with the variance, so the path encoder captures information beyond the current state when the variance is known.

\begin{figure}[htbp]
\centering
\resizebox{\textwidth}{!}{% Rough Bergomi Bermudan put: bounds of each method minus those of LogSig--NRDE, 95% intervals.
\begin{tikzpicture}
\pgfplotsset{forest/.style={paper, ymin=-0.6, ymax=6.6, y dir=reverse, ytick={0,...,6},
    ymajorgrids=false, xmajorgrids=true, axis lines*=box,
    xticklabel style={/pgf/number format/fixed}, scaled x ticks=false, clip=false},
  ci/.style={only marks, mark=*, mark size=1.7pt, color=cInk!80, mark options={fill=cInk!80},
    error bars/x dir=both, error bars/x explicit,
    error bars/error bar style={line width=0.8pt, color=cInk!80},
    error bars/error mark options={rotate=90, mark size=1.8pt, color=cInk!80}},
}
\begin{groupplot}[group style={group size=4 by 1, horizontal sep=0.3cm, yticklabels at=edge left},
  forest, width=4.1cm, height=5.0cm, yticklabels={{\textbf{\textcolor{cStop!80!black}{LogSig--NRDE (ours)}}},{Markov state $(t,S_t)$},{LSTM on LogSig},{Longstaff--Schwartz},{signature LS},{Markov state $(t,S_t,V_t)$},{LogSig--NRDE $+\,V_t$}}, yticklabel style={font=\scriptsize, text=cInk}]
\nextgroupplot[title={(a) $\Delta L_0$, $K=4$}, xmin=-0.08, xmax=0.08]
\fill[cStop!9] (axis cs:\pgfkeysvalueof{/pgfplots/xmin},\pgfkeysvalueof{/pgfplots/ymin}) rectangle (axis cs:0,\pgfkeysvalueof{/pgfplots/ymax});
\draw[cStop!85!black, line width=1.1pt] (axis cs:0,\pgfkeysvalueof{/pgfplots/ymin}) -- (axis cs:0,\pgfkeysvalueof{/pgfplots/ymax});
\addplot[only marks, mark=diamond*, mark size=3pt, color=cStop!85!black] coordinates {(0,0)};
\draw[cInk!40, densely dashed, line width=0.5pt] (axis cs:\pgfkeysvalueof{/pgfplots/xmin},4.5) -- (axis cs:\pgfkeysvalueof{/pgfplots/xmax},4.5);
\addplot[ci] coordinates {(-0.027,1) += (0.0220,0) -= (0.0220,0) (-0.027,2) += (0.0240,0) -= (0.0240,0) (-0.038,3) += (0.0220,0) -= (0.0200,0) (-0.041,4) += (0.0270,0) -= (0.0270,0) (0.001,5) += (0.0270,0) -= (0.0270,0) (0.021,6) += (0.0240,0) -= (0.0240,0)};
\nextgroupplot[title={(b) $\Delta U_0$, $K=4$}, xmin=-0.4, xmax=0.75]
\fill[cStop!9] (axis cs:0,\pgfkeysvalueof{/pgfplots/ymin}) rectangle (axis cs:\pgfkeysvalueof{/pgfplots/xmax},\pgfkeysvalueof{/pgfplots/ymax});
\draw[cStop!85!black, line width=1.1pt] (axis cs:0,\pgfkeysvalueof{/pgfplots/ymin}) -- (axis cs:0,\pgfkeysvalueof{/pgfplots/ymax});
\addplot[only marks, mark=diamond*, mark size=3pt, color=cStop!85!black] coordinates {(0,0)};
\draw[cInk!40, densely dashed, line width=0.5pt] (axis cs:\pgfkeysvalueof{/pgfplots/xmin},4.5) -- (axis cs:\pgfkeysvalueof{/pgfplots/xmax},4.5);
\addplot[ci] coordinates {(0.037,1) += (0.0140,0) -= (0.0150,0) (0.073,2) += (0.0710,0) -= (0.0480,0) (0.264,4) += (0.0420,0) -= (0.0400,0) (-0.026,5) += (0.0170,0) -= (0.0180,0) (-0.135,6) += (0.0160,0) -= (0.0160,0)};
\nextgroupplot[title={(c) $\Delta L_0$, $K=12$}, xmin=-0.08, xmax=0.1]
\fill[cStop!9] (axis cs:\pgfkeysvalueof{/pgfplots/xmin},\pgfkeysvalueof{/pgfplots/ymin}) rectangle (axis cs:0,\pgfkeysvalueof{/pgfplots/ymax});
\draw[cStop!85!black, line width=1.1pt] (axis cs:0,\pgfkeysvalueof{/pgfplots/ymin}) -- (axis cs:0,\pgfkeysvalueof{/pgfplots/ymax});
\addplot[only marks, mark=diamond*, mark size=3pt, color=cStop!85!black] coordinates {(0,0)};
\draw[cInk!40, densely dashed, line width=0.5pt] (axis cs:\pgfkeysvalueof{/pgfplots/xmin},4.5) -- (axis cs:\pgfkeysvalueof{/pgfplots/xmax},4.5);
\addplot[ci] coordinates {(-0.005,1) += (0.0360,0) -= (0.0360,0) (-0.008,2) += (0.0350,0) -= (0.0360,0) (0.025,3) += (0.0360,0) -= (0.0350,0) (-0.005,4) += (0.0430,0) -= (0.0400,0) (0.042,5) += (0.0420,0) -= (0.0410,0) (0.046,6) += (0.0280,0) -= (0.0290,0)};
\nextgroupplot[title={(d) $\Delta U_0$, $K=12$}, xmin=-0.4, xmax=0.75]
\fill[cStop!9] (axis cs:0,\pgfkeysvalueof{/pgfplots/ymin}) rectangle (axis cs:\pgfkeysvalueof{/pgfplots/xmax},\pgfkeysvalueof{/pgfplots/ymax});
\draw[cStop!85!black, line width=1.1pt] (axis cs:0,\pgfkeysvalueof{/pgfplots/ymin}) -- (axis cs:0,\pgfkeysvalueof{/pgfplots/ymax});
\addplot[only marks, mark=diamond*, mark size=3pt, color=cStop!85!black] coordinates {(0,0)};
\draw[cInk!40, densely dashed, line width=0.5pt] (axis cs:\pgfkeysvalueof{/pgfplots/xmin},4.5) -- (axis cs:\pgfkeysvalueof{/pgfplots/xmax},4.5);
\addplot[ci] coordinates {(0.099,1) += (0.0370,0) -= (0.0360,0) (-0.003,2) += (0.0250,0) -= (0.0220,0) (0.64,4) += (0.0500,0) -= (0.0500,0) (-0.13,5) += (0.0350,0) -= (0.0340,0) (-0.303,6) += (0.0340,0) -= (0.0350,0)};
\end{groupplot}
\end{tikzpicture}
}
\caption{Rough Bergomi Bermudan put ($\kappa=100$): bounds of each method minus those of the LogSig--NRDE model on the price path, with $95\pct$ interval; the shaded side favors the LogSig--NRDE model. Reference model: $[L_0,U_0]=[8.496,10.643]$ for $K=4$ and $[8.507,11.704]$ for $K=12$.}
\label{fig:rough}
\end{figure}

%----------------------------------------------------------------------
\subsection{Certificates and benchmarks}\label{sec:exp-benchmarks}
%----------------------------------------------------------------------

Every trained model yields a lower and an upper price bound (\cref{thm:primaldual}). Each reference value lies within two standard errors of the certified interval (\cref{tab:certificates}), and the interval narrows as the learning step decreases, as \cref{thm:primaldual}(b) predicts.

\begin{table}[htbp]
\centering\footnotesize
\caption{Prices and primal--dual certificates of the stop-gradient LogSig--NRDE model; standard errors in parentheses.}
\label{tab:certificates}
\setlength{\tabcolsep}{4pt}
\begin{tabular}{@{}l l c c c c c@{}}
\toprule
Contract & & Reference & $y_0$ & $L_0$ & $U_0$ & width [\%]\\
\midrule
\multirow{3}{*}{Bermudan put, $K=4$} & $\kappa=90$  & 2.4117 & 2.4146 & 2.4104 {\scriptsize(0.005)} & 2.4945 {\scriptsize(0.005)} & 3.5\\
                                     & $\kappa=100$ & 5.9566 & 5.9616 & 5.9508 {\scriptsize(0.010)} & 6.0683 {\scriptsize(0.006)} & 2.0\\
                                     & $\kappa=110$ & 11.7243 & 11.7321 & 11.7124 {\scriptsize(0.015)} & 11.8696 {\scriptsize(0.006)} & 1.3\\
Amerasian put, $K=4$ & & 3.7611 & 3.7729 & 3.7620 {\scriptsize(0.008)} & 3.7917 {\scriptsize(0.005)} & 0.8\\
Basket put, $d=5$, $K=4$ & & -- & 3.4885 & 3.4951 {\scriptsize(0.010)} & 3.5689 {\scriptsize(0.004)} & 2.1\\
\midrule
\multirow{3}{*}{Max-call, $d=2$, $K=9$} & $S_0=90$  & 8.075 & 8.090 & 8.099 {\scriptsize(0.027)} & 8.519 {\scriptsize(0.014)} & 5.2\\
                                        & $S_0=100$ & 13.902 & 13.921 & 13.934 {\scriptsize(0.034)} & 14.524 {\scriptsize(0.017)} & 4.2\\
                                        & $S_0=110$ & 21.345 & 21.364 & 21.372 {\scriptsize(0.039)} & 22.178 {\scriptsize(0.019)} & 3.8\\
\multirow{3}{*}{Max-call, $d=5$, $K=9$} & $S_0=90$  & 16.644 & 16.697 & 16.633 {\scriptsize(0.036)} & 17.357 {\scriptsize(0.020)} & 4.3\\
                                        & $S_0=100$ & 26.159 & 26.237 & 26.143 {\scriptsize(0.044)} & 27.116 {\scriptsize(0.023)} & 3.7\\
                                        & $S_0=110$ & 36.772 & 36.875 & 36.755 {\scriptsize(0.049)} & 37.956 {\scriptsize(0.026)} & 3.3\\
\bottomrule
\end{tabular}
\end{table}

The LogSig--NRDE policy is never dominated by Longstaff--Schwartz regression \cite{longstaff2001}, deep optimal stopping \cite{becker2019}, or signature Longstaff--Schwartz \cite{bayer2025primal} (\cref{fig:benchmarks}). Its advantage over regression grows with the dimension, and its dual bound is markedly tighter than the signature-based dual bound.

\begin{figure}[htbp]
\centering
\resizebox{\textwidth}{!}{% Benchmarks on common paths: bound of each benchmark minus that of LogSig--NRDE, 95% intervals.
\begin{tikzpicture}
\pgfplotsset{forest/.style={paper, ymin=-0.6, ymax=8.6, y dir=reverse, ytick={0,...,8},
    ymajorgrids=false, xmajorgrids=true, axis lines*=box,
    xticklabel style={/pgf/number format/fixed}, scaled x ticks=false, clip=false},
  ci/.style={only marks, mark=*, mark size=1.7pt, color=cInk!80, mark options={fill=cInk!80},
    error bars/x dir=both, error bars/x explicit,
    error bars/error bar style={line width=0.8pt, color=cInk!80},
    error bars/error mark options={rotate=90, mark size=1.8pt, color=cInk!80}},
}
\begin{groupplot}[group style={group size=2 by 1, horizontal sep=0.35cm, yticklabels at=edge left},
  forest, ymin=-0.6, width=6.2cm, height=6.2cm, yticklabels={{Bermudan put},{Amerasian put},{basket put, $d=5$},{max-call $d=2$, $S_0=90$},{max-call $d=2$, $S_0=100$},{max-call $d=2$, $S_0=110$},{max-call $d=5$, $S_0=90$},{max-call $d=5$, $S_0=100$},{max-call $d=5$, $S_0=110$}}, yticklabel style={font=\scriptsize, text=cInk}]
\nextgroupplot[title={(a) $\Delta L_0$}, xmin=-0.55, xmax=0.06,
  legend style={at={(1.03,-0.09)}, anchor=north, font=\scriptsize, legend columns=3, column sep=6pt}]
\fill[cStop!9] (axis cs:\pgfkeysvalueof{/pgfplots/xmin},\pgfkeysvalueof{/pgfplots/ymin}) rectangle (axis cs:0,\pgfkeysvalueof{/pgfplots/ymax});
\draw[cStop!85!black, line width=1.1pt] (axis cs:0,\pgfkeysvalueof{/pgfplots/ymin}) -- (axis cs:0,\pgfkeysvalueof{/pgfplots/ymax});
\addplot[ci, color=cFull!90!black, mark=*, mark options={fill=cFull!90!black}, error bars/error bar style={color=cFull!90!black}, error bars/error mark options={rotate=90, mark size=1.5pt, color=cFull!90!black}] coordinates {(-0.004,-0.24) += (0.0080,0) -= (0.0080,0) (-0.002,0.76) += (0.0020,0) -= (0.0020,0) (0.0,1.76) += (0.0060,0) -= (0.0060,0) (-0.046,2.76) += (0.0200,0) -= (0.0210,0) (-0.044,3.76) += (0.0270,0) -= (0.0280,0) (-0.064,4.76) += (0.0310,0) -= (0.0280,0) (-0.144,5.76) += (0.0390,0) -= (0.0440,0) (-0.191,6.76) += (0.0520,0) -= (0.0490,0) (-0.212,7.76) += (0.0610,0) -= (0.0570,0)};
\addlegendentry{Longstaff--Schwartz}
\addplot[ci, color=cEma!90!black, mark=square*, mark options={fill=cEma!90!black}, error bars/error bar style={color=cEma!90!black}, error bars/error mark options={rotate=90, mark size=1.5pt, color=cEma!90!black}] coordinates {(-0.065,0.0) += (0.0510,0) -= (0.0730,0) (-0.074,1.0) += (0.0060,0) -= (0.0060,0) (-0.065,2.0) += (0.0110,0) -= (0.0120,0)};
\addlegendentry{deep opt.\ stopping}
\addplot[ci, color=cLayer!90!black, mark=triangle*, mark options={fill=cLayer!90!black}, error bars/error bar style={color=cLayer!90!black}, error bars/error mark options={rotate=90, mark size=1.5pt, color=cLayer!90!black}] coordinates {(-0.009,0.24) += (0.0120,0) -= (0.0110,0) (-0.001,1.24) += (0.0020,0) -= (0.0010,0) (-0.059,2.24) += (0.0140,0) -= (0.0150,0) (-0.271,3.24) += (0.0350,0) -= (0.0320,0) (-0.197,4.24) += (0.0400,0) -= (0.0420,0) (-0.214,5.24) += (0.0520,0) -= (0.0540,0) (-0.403,6.24) += (0.0570,0) -= (0.0490,0) (-0.385,7.24) += (0.0590,0) -= (0.0510,0) (-0.437,8.24) += (0.0660,0) -= (0.0570,0)};
\addlegendentry{signature LS}
\node[font=\tiny, anchor=north west, inner sep=2pt] at (axis cs:0,-0.55) {\textcolor{cStop!80!black}{ours}};
\nextgroupplot[title={(b) $\Delta U_0$, signature LS}, xmin=0, xmax=3.9,
  xbar, bar width=6pt, bar shift=0pt, nodes near coords, point meta=rawx,
  nodes near coords style={font=\tiny, text=cInk, anchor=west, /pgf/number format/.cd, fixed, zerofill, precision=2}]
\fill[cStop!9] (rel axis cs:0,0) rectangle (rel axis cs:1,1);
\addplot[fill=cLayer!70, draw=cLayer!90!black] coordinates {(0.53,0) (0.22,1) (0.17,2) (1.48,3) (1.8,4) (2.03,5) (2.77,6) (3.01,7) (3.26,8)};
\end{groupplot}
\end{tikzpicture}
}
\caption{Benchmarks on common paths: lower bound of each benchmark and upper bound of signature LS minus those of the LogSig--NRDE model, with $95\pct$ interval; the shaded side favors the LogSig--NRDE model. Deep optimal stopping was not run on the max-call.}
\label{fig:benchmarks}
\end{figure}

From one to ten assets, the certified interval of the basket put barely widens, and the geometric basket call retains its accuracy (\cref{fig:dim}). The model size grows with the dimension of the local log-signature, while the training time stays between $290$ and $920$\,s across all dimensions (\cref{fig:dim}c).

\begin{figure}[htbp]
\centering
\resizebox{\textwidth}{!}{% Dimension d = 1, 2, 5, 10 (M = 2): certified interval width of the
% Bermudan basket put, price error of the geometric basket call, size.
\begin{tikzpicture}
\begin{groupplot}[group style={group size=3 by 1, horizontal sep=1.5cm},
                  paper, width=5.6cm, height=5.0cm,
                  xmode=log, xmin=0.8, xmax=12.5, xtick={1,2,5,10},
                  xticklabels={$1$,$2$,$5$,$10$}, xlabel={dimension $d$}]
\nextgroupplot[title={(a) Accuracy}, ylabel={[\%]}, ymin=0, ymax=3,
               legend style={at={(0.97,0.36)}, anchor=east}]
\addplot[sStop] table[x=d, y=berm_relgap] {data/dim.dat};
\addplot[color=cEma, mark=diamond*, mark size=2.6pt] table[x=d, y=geo_err] {data/dim.dat};
\legend{put: $(U_0-L_0)/L_0$, call: $|y_0-V_0|/V_0$}

\nextgroupplot[title={(b) Model size}, ylabel={parameters [$\times10^{5}$]},
               ymin=0, ymax=5.2, ytick={0,1,...,5}, legend style={at={(0.03,0.97)}, anchor=north west}]
\addplot[sStop] table[x=d, y expr=\thisrow{berm_params}/1e5] {data/dim.dat};
\addplot[color=cEma, mark=diamond*, mark size=2.6pt] table[x=d, y expr=\thisrow{geo_params}/1e5] {data/dim.dat};
\legend{Bermudan put ($K=4$), geometric call ($K=1$)}

\nextgroupplot[title={(c) Training time}, ylabel={time [s]},
               ymin=0, ymax=1000, legend style={at={(0.03,0.97)}, anchor=north west}]
\addplot[sStop] table[x=d, y=berm_time] {data/dim.dat};
\addplot[color=cEma, mark=diamond*, mark size=2.6pt] table[x=d, y=geo_time] {data/dim.dat};
\legend{Bermudan put, geometric call}
\end{groupplot}
\end{tikzpicture}
}
\caption{Dimension scaling at $M=2$. (a) Relative width of the certified interval (Bermudan basket put, $K=4$) and relative price error (geometric basket call). (b) Number of parameters. (c) Training time.}
\label{fig:dim}
\end{figure}
